\textbf{Phần 1. Trắc nghiệm lựa chọn}
\setcounterpageref{ex}{0}

\begin{ex} %[2H3N16]
Lít là đơn vị dùng để đo
\choice
{Khối lượng của vật}
{Độ dài đoạn thẳng}
{\True Dung tích (sức chứa) của chất lỏng}
{Thời gian trong ngày}
\loigiai{
Lít dùng để đo dung tích chất lỏng.\\
\textbf{Đáp án của câu hỏi là C}
}
\end{ex}

\begin{ex} %[2H3N16]
Lít được viết tắt là
\choice
{$\text{kg}$}
{$\text{cm}$}
{$\text{m}$}
{\True $\text{l}$}
\loigiai{
Lít viết tắt là $\text{l}$.\\
\textbf{Đáp án của câu hỏi là D}
}
\end{ex}

\begin{ex} %[2H3H16]
Tính: $15\text{ l} + 8\text{ l} = \dots$
\choice
{$22\text{ l}$}
{$24\text{ l}$}
{\True $23\text{ l}$}
{$25\text{ l}$}
\loigiai{
$15 + 8 = 23\text{ l}$.\\
\textbf{Đáp án của câu hỏi là C}
}
\end{ex}

\begin{ex} %[2H3H16]
Tính: $38\text{ l} - 15\text{ l} = \dots$
\choice
{$22\text{ l}$}
{$24\text{ l}$}
{$25\text{ l}$}
{\True $23\text{ l}$}
\loigiai{
$38 - 15 = 23\text{ l}$.\\
\textbf{Đáp án của câu hỏi là D}
}
\end{ex}

\begin{ex} %[2H3H16]
Thùng thứ nhất chứa $20\text{ l}$ nước mắm, thùng thứ hai chứa $15\text{ l}$ nước mắm. Cả hai thùng chứa số lít nước mắm là
\choice
{\True $35\text{ l}$}
{$30\text{ l}$}
{$40\text{ l}$}
{$45\text{ l}$}
\loigiai{
$20 + 15 = 35\text{ l}$.\\
\textbf{Đáp án của câu hỏi là A}
}
\end{ex}

\begin{ex} %[2H3H16]
Can to đựng $18\text{ l}$ dầu, can nhỏ đựng $9\text{ l}$ dầu. Can to đựng nhiều hơn can nhỏ số lít dầu là
\choice
{$8\text{ l}$}
{\True $9\text{ l}$}
{$10\text{ l}$}
{$27\text{ l}$}
\loigiai{
$18 - 9 = 9\text{ l}$.\\
\textbf{Đáp án của câu hỏi là B}
}
\end{ex}

\begin{ex} %[2H3V16]
Một can chứa $10\text{ l}$ nước, người ta rót ra các chai loại $2\text{ l}$. Rót được bao nhiêu chai?
\choice
{\True $5$ chai}
{$4$ chai}
{$6$ chai}
{$7$ chai}
\loigiai{
$10 : 2 = 5$ chai.\\
\textbf{Đáp án của câu hỏi là A}
}
\end{ex}

\begin{ex} %[2H3V16]
Điền dấu thích hợp: $24\text{ l} + 6\text{ l} \dots 32\text{ l}$.
\choice
{$>$}
{\True $<$}
{$=$}
{Không so sánh được}
\loigiai{
$24 + 6 = 30\text{ l} < 32\text{ l}$.\\
\textbf{Đáp án của câu hỏi là B}
}
\end{ex}

\begin{ex} %[2H3V16]
Tính giá trị biểu thức: $35\text{ l} + 15\text{ l} - 20\text{ l}$.
\choice
{$25\text{ l}$}
{$35\text{ l}$}
{\True $30\text{ l}$}
{$40\text{ l}$}
\loigiai{
$50 - 20 = 30\text{ l}$.\\
\textbf{Đáp án của câu hỏi là C}
}
\end{ex}

\begin{ex} %[2H3V16]
Một bể cá chứa $45\text{ l}$ nước, sau khi múc bớt $15\text{ l}$ nước thì trong bể còn lại bao nhiêu lít nước?
\choice
{$25\text{ l}$}
{\True $30\text{ l}$}
{$35\text{ l}$}
{$40\text{ l}$}
\loigiai{
$45 - 15 = 30\text{ l}$.\\
\textbf{Đáp án của câu hỏi là B}
}
\end{ex}

\begin{ex} %[2H3V16]
Bình đựng $5\text{ l}$ mật ong, can đựng $15\text{ l}$ mật ong. Hỏi can đựng gấp mấy lần bình?
\choice
{$2$ lần}
{$4$ lần}
{$5$ lần}
{\True $3$ lần}
\loigiai{
$15 : 5 = 3$ lần.\\
\textbf{Đáp án của câu hỏi là D}
}
\end{ex}

\begin{ex} %[2H3C16]
Có một can $10\text{ l}$ đựng đầy nước và hai can rỗng loại $3\text{ l}$ và $5\text{ l}$. Muốn lấy đúng $2\text{ l}$ nước, ta có thể
\choice
{\True Rót đầy can $5\text{ l}$ rồi từ can $5\text{ l}$ rót đầy can $3\text{ l}$, lượng nước còn lại trong can $5\text{ l}$ là $2\text{ l}$}
{Rót đầy can $3\text{ l}$}
{Rót nửa can $5\text{ l}$}
{Rót nửa can $10\text{ l}$}
\loigiai{
Lấy $5\text{ l} - 3\text{ l} = 2\text{ l}$.\\
\textbf{Đáp án của câu hỏi là A}
}
\end{ex}

\textbf{Phần 2. Trắc nghiệm đúng sai}
\setcounterpageref{ex}{0}

\begin{ex} %[2H3H16]
Xét tính đúng sai của các phép tính với đơn vị lít:
\choiceTF
{\True $18\text{ l} + 12\text{ l} = 30\text{ l}$.}
{\True $45\text{ l} - 25\text{ l} = 20\text{ l}$.}
{$10\text{ l} + 5 = 15\text{ l}$.}
{\True Khi cộng trừ dung tích phải viết kèm đơn vị $\text{l}$.}
\loigiai{
\begin{itemize}
    \item $18 + 12 = 30\text{ l}$: Đúng.
    \item $45 - 25 = 20\text{ l}$: Đúng.
    \item Ý 3 thiếu đơn vị ở số hạng thứ 2. Sai.
    \item Cần kèm đơn vị đo: Đúng.
\end{itemize}
\begin{itemize}
    \item a) \textbf{Đúng}.
    \item b) \textbf{Đúng}.
    \item c) \textbf{Sai}.
    \item d) \textbf{Đúng}.
\end{itemize}
}
\end{ex}

\begin{ex} %[2H3V16]
Bình nước có $16\text{ l}$, người ta dùng hết $7\text{ l}$. Xét tính đúng sai:
\choiceTF
{\True Lượng nước còn lại là $9\text{ l}$ ($16 - 7 = 9$).}
{\True Lượng nước ban đầu nhiều hơn lượng nước đã dùng là $9\text{ l}$.}
{Nếu đổ thêm $4\text{ l}$ nước vào bình thì trong bình có $15\text{ l}$.}
{\True Khi đó trong bình có $9 + 4 = 13\text{ l}$ nước.}
\loigiai{
\begin{itemize}
    \item Còn lại $9\text{ l}$: Đúng.
    \item $16 - 7 = 9\text{ l}$: Đúng.
    \item $9 + 4 = 13 \neq 15$. Ý 3 Sai, Ý 4 Đúng.
\end{itemize}
\begin{itemize}
    \item a) \textbf{Đúng}.
    \item b) \textbf{Đúng}.
    \item c) \textbf{Sai}.
    \item d) \textbf{Đúng}.
\end{itemize}
}
\end{ex}

\textbf{Phần 3. Trả lời ngắn}
\setcounterpageref{ex}{0}

\begin{ex} %[2H3H16]
Tính: $19\text{ l} + 8\text{ l}$ theo lít.
\shortans{27}
\loigiai{
\textbf{Đáp số:} $27$\\
$19 + 8 = 27\text{ l}$.
}
\end{ex}

\begin{ex} %[2H3H16]
Tính: $40\text{ l} - 15\text{ l}$ theo lít.
\shortans{25}
\loigiai{
\textbf{Đáp số:} $25$\\
$40 - 15 = 25\text{ l}$.
}
\end{ex}

\begin{ex} %[2H3V16]
Can to có $25\text{ l}$ nước mắm, can nhỏ có $15\text{ l}$ nước mắm. Cả hai can có tất cả bao nhiêu lít nước mắm?
\shortans{40}
\loigiai{
\textbf{Đáp số:} $40$\\
$25 + 15 = 40\text{ l}$.
}
\end{ex}

\begin{ex} %[2H3V16]
Một thùng có $50\text{ l}$ dầu, đã bán đi $30\text{ l}$ dầu. Hỏi trong thùng còn lại bao nhiêu lít dầu?
\shortans{20}
\loigiai{
\textbf{Đáp số:} $20$\\
$50 - 30 = 20\text{ l}$.
}
\end{ex}

\begin{ex} %[2H3V16]
Tính giá trị biểu thức: $20\text{ l} + 30\text{ l} - 15\text{ l}$ theo lít.
\shortans{35}
\loigiai{
\textbf{Đáp số:} $35$\\
$50 - 15 = 35\text{ l}$.
}
\end{ex}

\begin{ex} %[2H3C16]
Thùng thứ nhất có $18\text{ l}$ dầu, thùng thứ hai có ít hơn thùng thứ nhất $6\text{ l}$ dầu. Hỏi cả hai thùng có tất cả bao nhiêu lít dầu?
\shortans{30}
\loigiai{
\textbf{Đáp số:} $30$\\
Thùng 2: $18 - 6 = 12\text{ l}$. Cả hai thùng: $18 + 12 = 30\text{ l}$.
}
\end{ex}
